\chapter{ปริภูมิพิกัด 3 มิติ และการติดตามการเคลื่อนที่ 6 องศาอิสระ}
\label{chap:ch02}
\index{ปริภูมิสามมิติ (3D Space)}
\index{การติดตาม 6DoF (6-DoF Tracking)}
\index{ควอเทอร์นียัน (Quaternion)}

% ==========================================================
% แผนบริหารการสอนประจำบทที่ 2
% ==========================================================
\section*{แผนบริหารการสอนประจำบทที่ 2}
\addcontentsline{toc}{section}{แผนบริหารการสอนประจำบทที่ 2}

\noindent\textbf{หัวข้อเนื้อหาประจำบท}
\begin{enumerate}
    \item ระบบพิกัดสามมิติและปริภูมิอ้างอิงเชิงพื้นที่ในมาตรฐาน OpenXR
    \item การแปลงพิกัดวัตถุแข็งเกร็งด้วยเวกเตอร์การเลื่อนตำแหน่งและเมทริกซ์การหมุน
    \item สถาปัตยกรรมการตรวจจับและติดตามตำแหน่ง 6 องศาอิสระ (6DoF)
\end{enumerate}

\noindent\textbf{วัตถุประสงค์เชิงพฤติกรรม}
\begin{enumerate}
    \item จำแนกคุณลักษณะและความแตกต่างของ View, Local, และ Stage Reference Spaces ได้
    \item คำนวณการแปลงตำแหน่งและการหมุนในปริภูมิสามมิติด้วยควอเทอร์เนียนได้ถูกต้อง
    \item อธิบายหลักการหลอมรวมสัญญาณเซนเซอร์สำหรับการติดตามแบบ 6DoF ได้
    \item ประยุกต์ใช้การแปลงพิกัดในการวางตำแหน่งเครื่องมือทดลองเสมือนได้อย่างแม่นยำ
\end{enumerate}

\newpage

\noindent\textbf{กิจกรรมการเรียนการสอน}
\begin{enumerate}
    \item การบรรยายหลักการทางเรขาคณิตสามมิติและการสาธิตการหมุนด้วยควอเทอร์เนียน
    \item กิจกรรมการคำนวณเวกเตอร์ตำแหน่งและการแปลงพิกัดของอุปกรณ์ทดลองเสมือน
    \item การฝึกเขียนโปรแกรมกำหนดปริภูมิอ้างอิงและอ่านค่าพิกัด 6DoF ผ่าน OpenXR
    \item การทดลองจำลองปัญหาการติดขัดของแกนหมุน (Gimbal Lock) และการแก้ไขด้วยควอเทอร์เนียน
\end{enumerate}

\noindent\textbf{สื่อการเรียนการสอน}
\begin{enumerate}
    \item สไลด์บรรยายอิเล็กทรอนิกส์และเอกสารประกอบการสอนประจำบทที่ 2
    \item แบบจำลองเชิงภาพสามมิติแสดงปริภูมิพิกัดและมุมออยเลอร์
    \item โค้ดตัวอย่างการสร้าง Reference Space ในภาษา Python และ C++
    \item ระบบจำลองเซนเซอร์ติดตามตำแหน่งเสมือน (OpenXR Mock Tracking)
\end{enumerate}

\noindent\textbf{การประเมินผล}
\begin{enumerate}
    \item การตรวจแบบฝึกหัดคำนวณการคูณควอเทอร์เนียนและการแปลงพิกัดสามมิติ
    \item การประเมินผลงานโค้ดการสร้าง Spatial Anchor สำหรับอุปกรณ์ทดลองเสมือน
    \item การทดสอบความเข้าใจย่อยประจำสัปดาห์
\end{enumerate}

\newpage

\section{ปริภูมิอ้างอิงเชิงพื้นที่ใน OpenXR}

ในระบบประมวลผลเชิงพื้นที่ วัตถุเสมือนทุกชิ้นต้องมีจุดอ้างอิงตำแหน่งในปริภูมิสามมิติ มาตรฐาน OpenXR กำหนดประเภทของปริภูมิอ้างอิง (Reference Space Types) ไว้อย่างชัดเจนเพื่อรองรับสถานการณ์การใช้งานที่แตกต่างกัน ดังแสดงในตารางที่ 2.1

\begin{table}[htbp]
\centering
\caption{คุณลักษณะของปริภูมิอ้างอิงหลักในมาตรฐาน OpenXR}
\label{tab:openxr_spaces}
\begin{tabularx}{\textwidth}{lYY}
\toprule
\textbf{ชื่อปริภูมิอ้างอิง} & \textbf{จุดกำเนิดและทิศทางอ้างอิง} & \textbf{กรณีการใช้งานในสะเต็มศึกษา} \\
\midrule
\texttt{VIEW} & จุดกึ่งกลางระหว่างดวงตาทั้งสองข้างของกล้อง & การแสดงผล HUD หรือเมนูที่ติดตามสายตาผู้เรียน \\
\texttt{LOCAL} & จุดเริ่มต้นตำแหน่งของศีรษะขณะเริ่มทำงาน & การทดลองในท่านั่ง โต๊ะทดลองหน้าผู้เรียน \\
\texttt{STAGE} & กึ่งกลางพื้นที่ห้องบนพื้นราบจริง (Room-scale) & การเดินสำรวจห้องปฏิบัติการเสมือนจริงขนาดใหญ่ \\
\bottomrule
\end{tabularx}
\end{table}

\begin{stemconceptbox}{ปริภูมิพิกัดมือขวา (Right-Handed Coordinate System)}
มาตรฐาน OpenXR กำหนดให้ใช้ระบบพิกัดมือขวาในหน่วยเมตร (SI Unit) โดยแกน $+X$ ชี้ไปทางขวา แกน $+Y$ ชี้ขึ้นด้านบนในแนวดิ่ง และแกน $+Z$ ชี้ออกมาทางด้านหลัง (ซึ่งทำให้แกน $-Z$ ชี้ตรงไปข้างหน้าในทิศทางการมอง) ความสม่ำเสมอของหน่วยวัดทางฟิสิกส์นี้ช่วยให้การคำนวณสมการการเคลื่อนที่และการจำลองแรงทำได้อย่างตรงไปตรงมา
\end{stemconceptbox}

\section{การแปลงทางเรขาคณิตและควอเทอร์เนียน}

การระบุตำแหน่งและการวางตัวของวัตถุเสมือนในปริภูมิสามมิติประกอบด้วยเวกเตอร์การเลื่อนตำแหน่ง $\mathbf{p} = (x, y, z)^{\top}$ และการหมุน ซึ่งสามารถแทนได้ด้วยควอเทอร์เนียน $\mathbf{q}$

\subsection{นิยามของควอเทอร์เนียน}

ควอเทอร์เนียนเป็นโครงสร้างพีชคณิต 4 มิติที่ประกอบด้วยส่วนสเกลาร์และส่วนเวกเตอร์
\begin{equation}
\mathbf{q} = w + x\mathbf{i} + y\mathbf{j} + z\mathbf{k} = (w, \mathbf{v})
\end{equation}
โดยที่หน่วยจินตภาพสอดคล้องกับกฎของแฮมิลตัน
\begin{equation}
\mathbf{i}^2 = \mathbf{j}^2 = \mathbf{k}^2 = \mathbf{i}\mathbf{j}\mathbf{k} = -1
\end{equation}

เมื่อต้องการหมุนเวกเตอร์ตำแหน่ง $\mathbf{p}$ รอบแกนหมุนหนึ่งหน่วย $\mathbf{u} = (u_x, u_y, u_z)$ เป็นมุม $\theta$ เราสร้างควอเทอร์เนียนหนึ่งหน่วย (Unit Quaternion) ได้ดังนี้
\begin{equation}
\mathbf{q} = \cos\left(\frac{\theta}{2}\right) + \mathbf{u}\sin\left(\frac{\theta}{2}\right)
\end{equation}
และการหมุนเวกเตอร์ทำได้ผ่านการคูณแบบสังยุค (Conjugate Multiplication)
\begin{equation}
\mathbf{p}' = \mathbf{q} \mathbf{p} \mathbf{q}^*
\end{equation}
โดยที่ $\mathbf{q}^* = (w, -\mathbf{v})$ คือควอเทอร์เนียนสังยุค

\begin{figure}[htbp]
\centering
\begin{tikzpicture}[scale=1.2]
    % 3D Coordinate System
    \draw[-{Stealth[scale=1.2]}, line width=1.5pt, openxrTeal] (0,0,0) -- (4,0,0) node[anchor=north east]{$+X$ (ขวา)};
    \draw[-{Stealth[scale=1.2]}, line width=1.5pt, emeraldActive] (0,0,0) -- (0,3.5,0) node[anchor=north west]{$+Y$ (บน)};
    \draw[-{Stealth[scale=1.2]}, line width=1.5pt, spatialBlue] (0,0,0) -- (0,0,4) node[anchor=south east]{$+Z$ (หลัง)};

    % Rotation Axis u
    \draw[-{Stealth[scale=1.3]}, line width=1.8pt, amberGlow] (0,0,0) -- (2.5,3.0,1.5) node[above right]{\textbf{แกนหมุน $\mathbf{u}$}};

    % Arc of rotation theta
    \draw[->, line width=1.4pt, cyberPurple] (2.2, 1.8, 1.3) arc (20:200:0.8cm) node[midway, above]{\textbf{มุมหมุน $\theta$}};

    % Apparatus placement point
    \fill[darkNavy] (2.0, 1.2, 0.5) circle (3pt) node[below right]{\textbf{อุปกรณ์ทดลอง $\mathbf{p}'$}};
\end{tikzpicture}
\caption{การหมุนตำแหน่งของอุปกรณ์ทดลองในปริภูมิสามมิติด้วยแกนหมุนหนึ่งหน่วยและมุมหมุน}
\label{fig:quaternion_rotation}
\end{figure}

\begin{workedexamplebox}{การคำนวณการหมุนตำแหน่งหัววัดเซนเซอร์ด้วยควอเทอร์เนียน}
กำหนดให้เซนเซอร์วัดสนามแม่เหล็กติดตั้งอยู่ที่ตำแหน่งเริ่มต้น $\mathbf{p} = (0.5, 0, 0)\text{ m}$ ในปริภูมิห้องปฏิบัติการ เมื่อผู้ใช้หมุนโต๊ะทดลองรอบแกน $+Y$ ในแนวตั้งเป็นมุม $\theta = 90^\circ$ ตามเข็มนาฬิกาเมื่อมองจากด้านบน จงหาพิกัดตำแหน่งใหม่ $\mathbf{p}'$ ของเซนเซอร์

\textbf{วิธีทำ:}
แกนหมุนคือแกน $+Y$ ดังนั้น $\mathbf{u} = (0, 1, 0)$ และมุมหมุนตามเข็มนาฬิกาคือ $\theta = -90^\circ$ หรือ $-\frac{\pi}{2}\text{ rad}$
คำนวณองค์ประกอบของควอเทอร์เนียนหนึ่งหน่วย
\begin{align}
w &= \cos(-45^\circ) = \frac{\sqrt{2}}{2} \\
\mathbf{v} &= (0, 1, 0) \sin(-45^\circ) = \left(0, -\frac{\sqrt{2}}{2}, 0\right)
\end{align}
ดังนั้น $\mathbf{q} = \left(\frac{\sqrt{2}}{2}, 0, -\frac{\sqrt{2}}{2}, 0\right)$ และ $\mathbf{q}^* = \left(\frac{\sqrt{2}}{2}, 0, \frac{\sqrt{2}}{2}, 0\right)$

เมื่อทำการคูณ $\mathbf{p}' = \mathbf{q} \mathbf{p} \mathbf{q}^*$ โดยมอง $\mathbf{p} = (0, 0.5, 0, 0)$ จะได้
\begin{equation}
\mathbf{p}' = (0, 0, 0, 0.5)\text{ m}
\end{equation}
นั่นคือ ตำแหน่งใหม่ของหัววัดเซนเซอร์จะย้ายไปอยู่ที่ $x = 0\text{ m}, y = 0\text{ m}, z = 0.5\text{ m}$ ซึ่งสอดคล้องกับการหมุนทางกายภาพจริง
\end{workedexamplebox}

\subsection{การแปลงระบบพิกัดคาร์ทีเซียนและทรงกลมในระบบติดตาม 3 มิติ}

ในการพัฒนาสภาพแวดล้อมเสมือนจริงและการติดตามวัตถุเชิงพื้นที่ การแปลงระหว่างระบบพิกัดคาร์ทีเซียน (Cartesian Coordinates) $(x, y, z)$ และระบบพิกัดทรงกลม (Spherical Coordinates) $(r, \theta, \phi)$ มีบทบาทสำคัญอย่างยิ่งต่อการตรวจจับทิศทางการชี้ของลำแสง (Pointer Ray Casting) จากมือผู้ใช้ ตลอดจนการจัดวางส่วนต่อประสานกับผู้ใช้แบบลอยตัวในลักษณะทรงกลม (Spherical HUD Menu) ที่รักษาระยะห่างคงที่รอบดวงตาหรือศีรษะ

ความสัมพันธ์ในการแปลงจากระบบพิกัดคาร์ทีเซียนสามมิติไปสู่ระบบพิกัดทรงกลมในมาตรฐานกราฟิกคอมพิวเตอร์ (แกน $+Y$ อยู่ในแนวดิ่ง) แสดงได้ดังสมการ
\begin{align}
r &= \sqrt{x^2 + y^2 + z^2} \label{eq:spherical_r} \\
\theta &= \arccos\left(\frac{y}{r}\right) \quad (0 \le \theta \le \pi) \label{eq:spherical_theta} \\
\phi &= \operatorname{atan2}(z, x) \quad (-\pi < \phi \le \pi) \label{eq:spherical_phi}
\end{align}
โดยที่ $r$ คือระยะห่างในแนวรัศมีจากจุดกำเนิด $\theta$ คือมุมขั้ว (Polar Angle หรือ Zenith Angle) วัดเทียบกับแนวดิ่งแกน $+Y$ และ $\phi$ คือมุมทิศทางในระนาบแนวนอน (Azimuthal Angle) ในระนาบ $XZ$

ในทางกลับกัน เมื่อต้องการแปลงจากพิกัดทรงกลมกลับสู่พิกัดคาร์ทีเซียนสำหรับการกำหนดตำแหน่งวัตถุในเอนจินกราฟิก สามารถคำนวณได้จาก
\begin{align}
x &= r \sin\theta \cos\phi \label{eq:cartesian_x} \\
y &= r \cos\theta \label{eq:cartesian_y} \\
z &= r \sin\theta \sin\phi \label{eq:cartesian_z}
\end{align}

\begin{figure}[htbp]
\centering
\begin{tikzpicture}[scale=1.1]
    % Axes
    \draw[-{Stealth[scale=1.1]}, line width=1.2pt, slateGrey] (0,0,0) -- (4.2,0,0) node[anchor=north east]{$+X$};
    \draw[-{Stealth[scale=1.1]}, line width=1.2pt, slateGrey] (0,0,0) -- (0,3.8,0) node[anchor=north west]{$+Y$};
    \draw[-{Stealth[scale=1.1]}, line width=1.2pt, slateGrey] (0,0,0) -- (0,0,3.8) node[anchor=south east]{$+Z$};

    % Target Vector P
    \coordinate (P) at (2.8, 2.5, 2.0);
    \coordinate (Pproj) at (2.8, 0, 2.0);

    % Projection lines
    \draw[dashed, slateGrey!80!white, line width=0.8pt] (P) -- (Pproj);
    \draw[dashed, slateGrey!80!white, line width=0.8pt] (0,0,0) -- (Pproj);
    \draw[dashed, slateGrey!80!white, line width=0.8pt] (Pproj) -- (2.8, 0, 0);
    \draw[dashed, slateGrey!80!white, line width=0.8pt] (Pproj) -- (0, 0, 2.0);

    % Ray r
    \draw[-{Stealth[scale=1.2]}, line width=2.0pt, openxrTeal] (0,0,0) -- (P) node[above right]{\textbf{$\mathbf{P}(x, y, z)$}};
    \node[midway, above left, color=openxrTeal] at (1.4, 1.4, 1.0) {$r$};

    % Angle theta
    \draw[->, line width=1.2pt, cyberPurple] (0,1.5,0) arc (90:50:1.5) node[midway, above right]{\textbf{$\theta$}};

    % Angle phi in horizontal plane
    \draw[->, line width=1.2pt, amberGlow] (1.2,0,0) arc (0:40:1.2) node[midway, below right]{\textbf{$\phi$}};

    % Origin label
    \fill[darkNavy] (0,0,0) circle (2.5pt) node[below left]{$\mathbf{O}$ (จุดเซนเซอร์)};
\end{tikzpicture}
\caption{การแปลงเรขาคณิตระหว่างระบบพิกัดคาร์ทีเซียนสามมิติและระบบพิกัดทรงกลมในสภาพแวดล้อม XR}
\label{fig:cartesian_spherical}
\end{figure}

\begin{aipromptbox}[การคำนวณควอเทอร์เนียนและการสอดแทรกตำแหน่งเชิงมุม (SLERP) ด้วย AI]
\textbf{วัตถุประสงค์:} สร้างฟังก์ชัน Python สำหรับการสอดแทรกตำแหน่งเชิงมุมอย่างราบรื่น (Spherical Linear Interpolation: SLERP) ระหว่างควอเทอร์เนียนสองค่า สำหรับใช้ในการปรับทิศทางการมองเห็นและหมุนเครื่องมือทดลองเสมือนในสภาพแวดล้อม XR\\
\textbf{โครงสร้าง CLEAR Prompt:}
\begin{itemize}[topsep=1pt, itemsep=1pt, leftmargin=1.2em]
    \item \textbf{Context (บริบท):} ระบบจำลองฟิสิกส์เชิงพื้นที่บน OpenXR ต้องการการหมุนโมเดล 3 มิติแบบนุ่มนวล ไม่กระตุก
    \item \textbf{Logic (ตรรกะ):} ใช้นิยามการสอดแทรกบนผิวทรงกลม 4 มิติ $\mathbf{q}(t) = \frac{\sin((1-t)\Omega)}{\sin\Omega}\mathbf{q}_0 + \frac{\sin(t\Omega)}{\sin\Omega}\mathbf{q}_1$ พร้อมจัดการกรณีดอทโปรดักต์ติดลบ (Shortest Path)
    \item \textbf{Expectation (ผลลัพธ์ที่ต้องการ):} ฟังก์ชัน Python คลาส \texttt{Quaternion} พร้อมเมท็อด \texttt{slerp(q1, q2, t)} ที่รองรับ NumPy array และมี Unit test ตรวจสอบความถูกต้องที่ $t=0, 0.5, 1.0$
    \item \textbf{Action \& Restriction (ข้อกำหนด):} ห้ามใช้ external library อื่นนอกจาก \texttt{numpy} และต้องนอร์มัลไลซ์ผลลัพธ์ให้เป็น Unit Quaternion เสมอ
\end{itemize}
\end{aipromptbox}

\begin{codebox}[การสอดแทรกการหมุนสามมิติด้วย SLERP ในภาษา Python]
import numpy as np

def quaternion_slerp(q0, q1, t):
    """
    คำนวณการสอดแทรกเชิงมุมทรงกลม (SLERP) ระหว่างควอเทอร์เนียน q0 และ q1
    พารามิเตอร์: q0, q1 ในรูปแบบ [x, y, z, w], t ในช่วง [0.0, 1.0]
    """
    q0 = np.array(q0, dtype=np.float64) / np.linalg.norm(q0)
    q1 = np.array(q1, dtype=np.float64) / np.linalg.norm(q1)
    
    # คำนวณ Dot Product ระหว่างสองควอเทอร์เนียน
    dot = np.dot(q0, q1)
    
    # หาก Dot Product ติดลบ ให้กลับเครื่องหมายเพื่อเลือกเส้นทางที่สั้นที่สุด
    if dot < 0.0:
        q1 = -q1
        dot = -dot
        
    # หากมุมชิดกันมาก (dot > 0.9995) ให้ใช้ Linear Interpolation แทนเพื่อป้องกันหารด้วยศูนย์
    if dot > 0.9995:
        result = q0 + t * (q1 - q0)
        return result / np.linalg.norm(result)
        
    theta_0 = np.arccos(np.clip(dot, -1.0, 1.0))
    theta = theta_0 * t
    q2 = q1 - q0 * dot
    q2 = q2 / np.linalg.norm(q2)
    
    return q0 * np.cos(theta) + q2 * np.sin(theta)

# ตัวอย่างการทดสอบ: หมุนจาก 0 องศา ไป 90 องศารอบแกน Y ที่ t = 0.5 (45 องศา)
q_start = [0.0, 0.0, 0.0, 1.0]
q_end = [0.0, np.sin(np.pi/4), 0.0, np.cos(np.pi/4)]
q_mid = quaternion_slerp(q_start, q_end, 0.5)
print(f"Interpolated Quaternion at t=0.5: {np.round(q_mid, 4)}")
\end{codebox}

\section{การติดตามการเคลื่อนที่ 6 องศาอิสระ}

การติดตามแบบ 6DoF (Six Degrees of Freedom) ประกอบด้วยการเคลื่อนที่เชิงเส้น 3 ทิศทาง ($X, Y, Z$) และการหมุนเชิงมุม 3 ทิศทาง (Yaw, Pitch, Roll) ในอุปกรณ์ XR สมัยใหม่ การติดตามอาศัยการหลอมรวมสัญญาณเซนเซอร์ (Sensor Fusion) ระหว่างหน่วยวัดความเฉื่อย (Inertial Measurement Unit: IMU) ที่ทำงานด้วยความถี่สูง ($>1000\text{ Hz}$) ร่วมกับกล้องติดตามทัศนวิสัยภายนอกสู่ภายใน (Inside-Out Optical Tracking) ที่ทำงานด้วยความถี่ $60\text{ Hz}$ เพื่อขจัดปัญหาการเลื่อนไถลสะสม (Drift Elimination)

\begin{codebox}[การสร้างปริภูมิอ้างอิง Stage ใน OpenXR ด้วยภาษา Python]
import openxr as xr

# กำหนดพารามิเตอร์การสร้าง Reference Space
create_info = xr.ReferenceSpaceCreateInfo(
    reference_space_type=xr.ReferenceSpaceType.STAGE,
    pose_in_reference_space=xr.Posef(
        orientation=xr.Quaternionf(0.0, 0.0, 0.0, 1.0),
        position=xr.Vector3f(0.0, 0.0, 0.0)
    )
)

# สร้าง Spatial Reference Space จากเซสชัน
stage_space = xr.create_reference_space(session, create_info)
print("Stage Space initialized successfully for Room-Scale STEM Lab.")
\end{codebox}

\subsection{ขั้นตอนวิธี Unprojection และการแมปพิกัดจอภาพ 2D สู่ปริภูมิ 3D}

เมื่อผสานการทำงานระหว่างการตรวจจับวัตถุหรือมือผ่านกล้อง 2 มิติ (เช่น การใช้เว็บแคมร่วมกับโมเดลการประมวลผลภาพ) เพื่อควบคุมและวางตำแหน่งเครื่องมือทดลองในปริภูมิ 3 มิติ จำเป็นต้องมีขั้นตอนวิธี \textbf{Camera Unprojection} เพื่อย้อนการฉายภาพจากพิกัดระนาบจอภาพ $(x_s, y_s)$ ในหน่วยพิกเซล กลับไปเป็นลำแสงเวกเตอร์สามมิติในปริภูมิโลก (World Space)

ขั้นตอนวิธีเริ่มต้นจากการแปลงพิกัดพิกเซลบนหน้าจอขนาดความกว้าง $W$ และความสูง $H$ ให้อยู่ในระบบพิกัดอุปกรณ์นอร์มัลไลซ์ (Normalized Device Coordinates: NDC) ในช่วง $[-1, 1]$
\begin{align}
x_{\text{ndc}} &= \frac{2 x_s}{W} - 1 \label{eq:ndc_x} \\
y_{\text{ndc}} &= 1 - \frac{2 y_s}{H} \label{eq:ndc_y}
\end{align}

จากนั้นนำเวกเตอร์พิกัด $\mathbf{v}_{\text{ndc}} = (x_{\text{ndc}}, y_{\text{ndc}}, z_{\text{ndc}}, 1)^{\top}$ โดยเลือก $z_{\text{ndc}} = 1.0$ (ที่ระนาบตัดภาพด้านไกล หรือ Far Clipping Plane) คูณด้วยเมทริกซ์ผกผันของการฉายภาพและการมองเห็น $\mathbf{M}_{\text{inv}} = (\mathbf{P} \mathbf{V})^{-1}$
\begin{equation}
\mathbf{v}_{\text{world}} = (\mathbf{P} \mathbf{V})^{-1} \begin{pmatrix} x_{\text{ndc}} \\ y_{\text{ndc}} \\ 1.0 \\ 1.0 \end{pmatrix}
\end{equation}
เมื่อแปลงกลับจากพิกัดเอกพันธุ์ (Homogeneous Coordinates) โดยการหารด้วยองค์ประกอบ $w$ จะได้พิกัดปลายทาง $\mathbf{P}_{\text{far}}$ และสามารถคำนวณเวกเตอร์ทิศทางของลำแสงหนึ่งหน่วย $\mathbf{d}_{\text{ray}}$ จากตำแหน่งกล้อง $\mathbf{C}_{\text{cam}}$ ได้ดังนี้
\begin{equation}
\mathbf{d}_{\text{ray}} = \frac{\mathbf{P}_{\text{far}} - \mathbf{C}_{\text{cam}}}{\|\mathbf{P}_{\text{far}} - \mathbf{C}_{\text{cam}}\|}
\end{equation}
วัตถุเสมือนที่ระยะความลึก $d_{\text{depth}}$ จากตำแหน่งกล้องจึงมีพิกัด 3 มิติในโลกจำลองคือ
\begin{equation}
\mathbf{P}_{\text{target}} = \mathbf{C}_{\text{cam}} + d_{\text{depth}} \mathbf{d}_{\text{ray}}
\end{equation}

\subsection{การจัดการหน่วยความจำแบบ Zero-GC Vector Pooling สำหรับ 60 FPS}

ในแอปพลิเคชันจำลองฟิสิกส์และการประมวลผลเชิงพื้นที่บนอุปกรณ์ XR หรือ WebXR ที่ต้องทำงานด้วยอัตราการเรนเดอร์ต่อเนื่อง 60 ถึง 120 เฟรมต่อวินาที ข้อผิดพลาดที่พบบ่อยในการเขียนโค้ดคือการสร้างอินสแตนซ์ของเวกเตอร์ใหม่ในทุกเฟรม เช่น \texttt{new Vector3()} ภายในฟังก์ชันการวนรอบเรนเดอร์ (Render Loop)

การสร้างอ็อบเจกต์ชั่วคราวอย่างต่อเนื่องจะส่งผลให้ตัวเก็บขยะหน่วยความจำ (Garbage Collector: GC) ต้องหยุดการทำงานของระบบเป็นระยะสั้นๆ (GC Pause หรือ Jank) ทำให้เฟรมเรตตกและผู้ใช้เกิดอาการเวียนศีรษะจากภาพกระตุก (Cyber-sickness)

แนวทางการแก้ไขตามมาตรฐานระดับมืออาชีพคือการใช้สถาปัตยกรรม \textbf{Zero-GC Vector Pooling} โดยการสร้างตัวแปรเวกเตอร์สำหรับทดเลขล่วงหน้า (Pre-allocated Scratch Vectors) ไว้นอกลูป แล้วนำกลับมาใช้ซ้ำ (Reuse) ตลอดการทำงาน ดังตัวอย่างโค้ดในตัวอย่างที่ 2.4

\begin{codebox}[การประยุกต์ใช้ Unprojection และ Zero-GC Vector Pool ใน JavaScript/Three.js]
// ประกาศ Pre-allocated Scratch Vectors ไว้นอก Render Loop เพื่อเลี่ยง GC Spikes
const _camPos = new THREE.Vector3();
const _vec = new THREE.Vector3();
const _targetPos = new THREE.Vector3();

function unprojectScreenToWorld(screenX, screenY, depth, camera, outPosition) {
    const width = window.innerWidth;
    const height = window.innerHeight;

    // 1. แปลงพิกเซลเป็น NDC (-1 ถึง +1)
    const ndcX = (screenX / width) * 2 - 1;
    const ndcY = -(screenY / height) * 2 + 1;

    // 2. กำหนดค่าใน Scratch Vector แทนการ new Vector3
    _vec.set(ndcX, ndcY, 0.5);

    // 3. Unproject ด้วย Matrix ผกผันของ Camera
    _vec.unproject(camera);

    // 4. คำนวณเวกเตอร์ทิศทางรังสีจากตำแหน่งกล้อง
    camera.getWorldPosition(_camPos);
    _vec.sub(_camPos).normalize();

    // 5. คำนวณตำแหน่ง 3D ที่ระยะ depth แล้วบันทึกลง outPosition (Zero-GC)
    outPosition.copy(_camPos).addScaledVector(_vec, depth);
    return outPosition;
}

// ใช้งานใน Render Loop:
// unprojectScreenToWorld(hand2DX, hand2DY, 0.8, camera, virtualBeaker.position);
\end{codebox}

\subsection{การวิเคราะห์ความคลาดเคลื่อนและการถดถอยเชิงเส้นในการแปลงพิกัด}

ในการทดลองสะเต็มศึกษา ข้อมูลพิกัดที่ได้จากการติดตามเซนเซอร์ย่อมมีความคลาดเคลื่อน (Noise) ปะปนอยู่เสมอ การตรวจสอบความสัมพันธ์เชิงเส้นระหว่างค่าพิกัดที่แท้จริง ($x$) กับค่าพิกัดที่อ่านได้จากระบบเซนเซอร์ ($y$) ทำได้ด้วยระเบียบวิธีถดถอยเชิงเส้นกำลังสองน้อยที่สุด (Least Squares Linear Regression) เพื่อหาความชัน $m$ และจุดตัดแกน $c$ ในสมการ $y = mx + c$
\begin{align}
m &= \frac{N \sum_{i=1}^N x_i y_i - \sum_{i=1}^N x_i \sum_{i=1}^N y_i}{N \sum_{i=1}^N x_i^2 - \left(\sum_{i=1}^N x_i\right)^2} \label{eq:linear_slope} \\
c &= \frac{\sum_{i=1}^N y_i - m \sum_{i=1}^N x_i}{N} \label{eq:linear_intercept}
\end{align}
เมื่อค่าสัมประสิทธิ์การตัดสินใจ $R^2 \to 1.0$ แสดงว่าระบบแปลงพิกัดมีความแม่นยำสูงและมีความถูกต้องเชิงเรขาคณิต ซึ่งข้อมูลการวิเคราะห์นี้สามารถนำไปใช้ปรับเทียบ (Calibrate) อุปกรณ์ทดลองเสมือนได้อย่างเป็นระบบ

\clearpage

\section*{สรุปท้ายบทที่ 2}
\addcontentsline{toc}{section}{สรุปท้ายบทที่ 2}

ปริภูมิอ้างอิงสามมิติและคณิตศาสตร์ของควอเทอร์เนียนเป็นเสาหลักสำคัญของการประมวลผลเชิงพื้นที่ในมาตรฐาน OpenXR การใช้หน่วยวัดสากลและระบบพิกัดมือขวาช่วยให้การเชื่อมโยงข้อมูลเซนเซอร์กับการสร้างอุปกรณ์ทดลองทางฟิสิกส์เป็นไปอย่างแม่นยำ การเข้าใจขั้นตอนวิธี Unprojection และการประยุกต์ใช้แนวคิด Zero-GC Vector Pooling เป็นหัวใจสำคัญในการคงประสิทธิภาพการทำงานที่ระดับ 60 ถึง 120 เฟรมต่อวินาที ป้องกันปัญหาการติดขัดของแกนหมุน และสร้างความพร้อมสำหรับการสร้างการปฏิสัมพันธ์ผ่านมือในบทต่อไป

\clearpage

\section*{แบบฝึกหัดท้ายบทที่ 2}
\addcontentsline{toc}{section}{แบบฝึกหัดท้ายบทที่ 2}

\begin{enumerate}
    \item จงเปรียบเทียบข้อดีและข้อเสียของการแทนการหมุนสามมิติด้วยมุมออยเลอร์กับควอเทอร์เนียน พร้อมอธิบายปรากฏการณ์ Gimbal Lock
    \item วัตถุมีตำแหน่งเริ่มต้นที่ $\mathbf{p} = (1, 2, -3)$ ในปริภูมิ Local จงแปลงพิกัดเมื่อวัตถุถูกเลื่อนด้วยเวกเตอร์ $\mathbf{t} = (-1, 0, 2)$ และหมุนรอบแกน $+Z$ เป็นมุม $180^\circ$
    \item อธิบายหลักการทำงานของ Inside-Out Optical Tracking ในการกำจัดความคลาดเคลื่อนสะสมของเซนเซอร์ IMU
    \item จุดบนหน้าจอความละเอียด $1920 \times 1080$ พิกเซล อยู่ที่พิกัด $(960, 540)$ จงคำนวณพิกัด Normalized Device Coordinates (NDC) ของจุดดังกล่าว
    \item อธิบายสาเหตุที่การเรียกคำสั่งสร้างอ็อบเจกต์เวกเตอร์ใหม่ในทุกเฟรม ส่งผลเสียต่อประสบการณ์ใช้งานในระบบเสมือนจริง และอธิบายวิธีแก้ปัญหาด้วย Zero-GC Vector Pooling
\end{enumerate}
